% CLASSIFY: L1
\documentclass[11pt]{article}
\usepackage{amsmath,amssymb,amsthm}
\usepackage[margin=1in]{geometry}
\usepackage{booktabs}
\usepackage{array}
\usepackage[colorlinks=true,linkcolor=blue,citecolor=blue,urlcolor=blue]{hyperref}

\newtheorem{theorem}{Theorem}[section]
\newtheorem{proposition}[theorem]{Proposition}
\newtheorem{lemma}[theorem]{Lemma}
\newtheorem{corollary}[theorem]{Corollary}
\theoremstyle{definition}
\newtheorem{definition}[theorem]{Definition}
\newtheorem{example}[theorem]{Example}
\theoremstyle{remark}
\newtheorem{remark}[theorem]{Remark}

\newcommand{\rid}[1]{\texttt{#1}}
\newcommand{\lean}[1]{\texttt{#1}}

\title{A Machine-Checked Kernel for Federated Automated Adjudication:\\
Credibility Lattice Completion, Certificate Lifecycle Invariants,\\
Certified Numerical Anchors, and a Parameterized Rice Template}
\author{PIVOT-01 Federation (machine-generated; no human author claimed;\\
AI-use disclosure attached as supplementary material)}
\date{October 11, 2026 (v2; supersedes the registration-style v1, identifier 2610.00183v1)}

\begin{document}
\maketitle

\begin{abstract}
We present a formally specified and machine-checked kernel for a federated
automated adjudication system, in which independent computational lines judge
mathematical and operational claims and deposit their verdicts, evidence, and
certificates in a shared ledger. The kernel comprises five components.
(i)~A three-rung, three-track credibility lattice obtained as the free
completion of a legacy five-element poset: the completion is justified by an
exhaustive enumeration over all $11^{3}=1331$ ordered triples and is fully
re-formalized in Lean~4, where the equational lattice laws, the order axioms,
and the extremality properties of join and meet are all proved by the
\texttt{decide} procedure over the finite carrier.
(ii)~A five-state certificate lifecycle machine with nine legal transitions
and machine-checked invariants, likewise re-formalized in Lean~4; the
formalization caught a specification flaw in one invariant (the originally
stated grant-gating condition is false), and we report the corrected,
machine-proved form.
(iii)~Certified numerical anchors: uniqueness of the gauged fixed point of
circulant entropic-optimal-transport (Sinkhorn) systems, established by a
Krawczyk interval operator with outward rounding across a six-instance family
($k\in\{6,10\}$, $\varepsilon\in\{1,1/2,1/5\}$, later extended to ten
instances), together with a second anchor family of five auction-algorithm
assignment instances certified by interval dual brackets. We state plainly
that the certification of these numerical anchors is carried by the Python
interval-arithmetic layer; the attempted port to the Lean~4 interval library
\emph{leancert} is blocked by structural capability boundaries of that
library (absence of the logarithm from its supported expression grammar and
infeasibility of its determinant computation at the required dimensions),
and the Lean files are formal statements, not machine-verified certificates.
(iv)~A parameterized Lean~4 template for Rice-type undecidability, consisting
of a bridge theorem, an extensionality lemma, an instantiation generator, and
discharged instances of increasing difficulty.
(v)~A minimal trusted computing base result: the soundness of an interval
certificate checker (accept implies correctness), itself machine-proved in
Lean~4. All Lean artifacts were verified on cloud infrastructure (Axiom Axle,
environment \texttt{lean-4.28.0}) under the strict \texttt{verify\_proof}
channel, with axiom audits restricted to \texttt{propext},
\texttt{Classical.choice}, and \texttt{Quot.sound}, and no \texttt{sorryAx}.
Empirical fuzzing over three certificate families observed zero false
accepts; we report rule-of-three 95\% upper bounds on the false-accept rate
and explicitly do not claim a guarantee.
\end{abstract}

\tableofcontents

\section{Introduction}\label{sec:intro}

\subsection{The problem: adjudication among untrusted computational lines}

Consider a federation of autonomous computational agents---``lines'' in the
terminology of our federation---that make claims about mathematical and
operational facts: that a nonlinear system has a unique solution in a box,
that a discrete optimization heuristic came within a certified bracket of the
optimum, that a proposed inference rule preserves truth. The lines do not
trust one another, and there is no central authority whose judgment is taken
on faith. Any claim that is to be used downstream must therefore carry
\emph{machine-checkable evidence}, and the rules by which evidence is
credited, promoted, maintained, downgraded, and revoked must themselves be
small enough to audit and rigid enough to be machine-checked.

This is the federated automated adjudication problem. It is not primarily a
problem of proving one hard theorem; it is a problem of \emph{governance
infrastructure for proofs}: a credibility order over claims, a lifecycle over
certificates, numerical anchors whose acceptance cannot be fabricated by
construction, a formal explanation of why fully automatic adjudication over
unrestricted semantic properties is impossible, and a minimal trusted base in
which the accept-implies-correct pattern is proved rather than assumed.

A naive design would adjudicate every semantic property of every program or
certificate. Rice's theorem~\cite{rice1953} forbids this: any nontrivial
extensional property of partial recursive functions is undecidable. The
federation therefore adopts a \emph{domain-limited, three-valued}
architecture: adjudication objects are quadruples
$o=(\varphi, D, \pi, V)$ consisting of a proposition $\varphi$, an explicit
applicability domain $D$, a certificate $\pi$, and a checker $V$; the
adjudication map $J$ is total by construction, taking values in
$\{\mathrm{pass}, \mathrm{fail}, \mathrm{undecided}\}$, with non-convergence
within a round assigned the third value. Within this architecture, the
soundness interface is the single clause
\[
V(\varphi, D, \pi) = \mathrm{accept} \;\Longrightarrow\; \varphi \text{ holds on } D,
\]
backed at each layer by a classical theorem: outward-rounding inclusion for
interval evaluation, the Krawczyk theorem for existence and uniqueness, and
LP weak duality for optimality brackets.

\subsection{What this paper deposits}

This paper is the full paper-form account of the kernel previously deposited
in registration form (v1, identifier 2610.00183v1). Every component below is
accompanied by machine certificates that are pinned, by content and by commit
identifier, in the federation's public repository, and every Lean~4 artifact
carries a permanent cloud-verification request identifier under which the
verification can be re-queried. The five components are:

\begin{enumerate}
\item \textbf{Credibility lattice completion} (Section~\ref{sec:lattice}).
The legacy five-element credibility poset has exactly seven join gaps, found
by enumeration; its completion is an eleven-element structure (three rungs
$\times$ three tracks, plus formal top and bottom). The lattice laws are
verified exhaustively over all $11^{3}=1331$ ordered triples, and the entire
structure is re-formalized in Lean~4 with every law decided over the finite
carrier (certificate CERT-LATTICE-01\,/\,CERT-LATTICE-LEAN-01).
\item \textbf{Certificate lifecycle machine} (Section~\ref{sec:lifecycle}).
A five-state machine (\emph{candidate}, \emph{granted}, \emph{maintained},
\emph{demoted}, \emph{revoked}) with exactly nine legal and eleven illegal
distinct transitions, with machine-checked gating, terminality, and freezing
invariants (CERT-K4-01\,/\,CERT-K4-LEAN-01). The Lean formalization refuted
our original stronger statement of the grant-gating invariant; the corrected
invariant is the machine-proved one.
\item \textbf{Certified numerical anchors} (Section~\ref{sec:anchors}). A
six-instance (later ten-instance) family of gauged circulant Sinkhorn
fixed-point systems, each certified to contain exactly one real root by a
Krawczyk interval operator with outward rounding and negative controls; and a
five-instance family of auction-algorithm assignment problems certified by
interval dual brackets (CERT-CIRC-01 family, CERT-AUCTION-01). The
certification carrier is the Python interval layer; we report, as a first-class
result of the project, the structural boundary that blocks the Lean-side port
(FINDING-04).
\item \textbf{Parameterized Rice template} (Section~\ref{sec:rice}). A
Lean~4 template that reduces any nontrivial extensional predicate on codes to
Mathlib's formalized Rice theorem, together with an instantiation generator
and seven discharged instances of increasing witness difficulty
(T2A-BRIDGE-01, CERT-T2A-TEMPLATE-01).
\item \textbf{Minimal trusted base} (Section~\ref{sec:selfcheck}). A Lean~4
proof that the interval certificate checker is sound---acceptance implies the
semantic property---for the step algebra used by the numerical anchors
(CERT-SELFCHECK-01).
\end{enumerate}

\subsection{Contributions}

The contributions of this paper are infrastructural and methodological rather
than a new theorem of independent depth:
\begin{itemize}
\item a complete, pinned, and machine-checked specification of an
adjudication kernel, with every enumeration certificate reproducible from
pinned scripts and every Lean artifact carrying a cloud-verification request
identifier;
\item a worked example in which formalization \emph{caught} a specification
flaw in a stated invariant (Section~\ref{sec:lifecycle}), which we regard as
the strongest argument for the kernel's methodology;
\item an honest account of a failed port: the circulant family is certified
at the Python interval layer, and its Lean-side formalization is structurally
unverifiable under the pinned version of the target library; we document the
exact boundary (Section~\ref{sec:finding04});
\item a reproducibility and axiom-audit protocol (Section~\ref{sec:repro}),
including rule-of-three honesty bounds on all empirical zero-failure claims.
\end{itemize}

\paragraph{Provenance and authorship stance.}
This paper is machine-generated by the PIVOT-01 federation orchestrator as
part of an explicit experiment in machine-checked adjudication
infrastructure. No human author is claimed; the contributing account holder
does not claim human understanding of every derivation step. Good-faith
correctness belief rests on the exhaustiveness of the enumeration
certificates, the outward-rounding discipline of the interval certificates,
kernel-level and cloud-level verification of the Lean artifacts, and the
rule-of-three statistical framing of all empirical claims. An AI-use
disclosure accompanies the submission as supplementary material, and the
derivation transcripts are retained in the federation ledger.

\paragraph{Organization.}
Section~\ref{sec:related} situates the work. Sections~\ref{sec:lattice}
through~\ref{sec:selfcheck} present the five kernel components.
Section~\ref{sec:repro} gives the reproducibility and axiom-audit protocol
and the evidence-anchor table. Section~\ref{sec:limitations} states
limitations and honest gaps. Section~\ref{sec:conclusion} concludes.
Appendix~\ref{app:anchors} collects the pinned artifact index.

\section{Related Work}\label{sec:related}

\paragraph{Proof assistants and Mathlib.}
Our formalizations are written in Lean~4~\cite{lean4} and build on
Mathlib~\cite{mathlib}, whose library of formalized mathematics includes a
development of computability theory---partial recursive functions, codes, and
Rice's theorem in the form \lean{ComputablePred.rice\_2}---due to
Carneiro~\cite{carneiro2019}. Our Rice template
(Section~\ref{sec:rice}) is a thin, reusable layer on top of that foundation:
the hard direction is entirely Mathlib's, and our contribution is the
parameterization and the instance generator. Our trust posture toward the
proof assistant itself follows the spirit of work on the metatheory and
self-verification of Lean, such as Lean4Lean~\cite{lean4lean}; we do not
rely on such results, and instead minimize and audit what our artifacts
assume (Section~\ref{sec:repro}).

\paragraph{Interval arithmetic and Krawczyk certification.}
Interval arithmetic as a tool for rigorous numerical bounds goes back to
Moore~\cite{moore1966}. The Krawczyk operator~\cite{krawczyk1969} converts a
candidate root and a box into a certified uniqueness statement: if the
operator maps the box strictly into itself, the box contains exactly one
root. Our numerical anchors (Section~\ref{sec:anchors}) use exactly this
criterion, with rationalized centers, rationalized analytic-inverse
preconditioners, and outward rounding at every step, so that the acceptance
of a certificate is a theorem about real numbers rather than about floating
point. The accept-implies-correct shape of these certificates is what our
minimal trusted base (Section~\ref{sec:selfcheck}) formalizes in miniature.

\paragraph{Optimal transport and auction algorithms.}
Entropic regularization of optimal transport and its Sinkhorn iteration are
classical~\cite{sinkhorn1967,cuturi2013}. We do not certify the iteration;
we certify the \emph{fixed-point system} directly for circulant cost
matrices, where a closed-form solution exists and serves as ground truth for
negative controls. The auction algorithm for assignment problems is due to
Bertsekas~\cite{bertsekas1988}; our second anchor family certifies auction
output a posteriori via interval dual brackets (LP weak duality), with the
solver treated as an untrusted candidate generator.

\paragraph{Undecidability and the shape of the escape.}
Rice's theorem~\cite{rice1953} is the formal reason that the federation's
adjudication map is three-valued and domain-limited. We treat the theorem not
as background but as a component of the kernel: the template of
Section~\ref{sec:rice} makes the blocking result callable, so that any
proposed extension of the adjudication domain to a new semantic property must
first pass (or be blocked by) a discharged instance.

\paragraph{Registration platforms for machine-generated mathematics.}
This submission is itself an experiment in the registration of
machine-generated, machine-checked mathematical content: the authorship mode
is \emph{no-human-author-claimed}, all artifacts are pinned in a public
repository, and the derivation transcripts are retained. A predecessor of
this paper was desk-rejected in moderation as ``not a paper''; the present
version is a deliberate test of whether a fully paper-form treatment of the
same evidence base meets the venue's expectations, and the outcome of that
test is itself data for the federation's platform-adaptation ledger.

\paragraph{Statistical honesty for zero-failure observations.}
When $n$ independent trials observe zero failures, the rule of
three~\cite{hanley1983} gives $3/n$ as an approximate 95\% upper confidence
bound on the failure rate. All empirical zero-false-accept claims in this
paper are stated with such bounds and with the explicit wording ``no false
accept observed'', never as guarantees.

\section{Credibility Lattice Completion}\label{sec:lattice}

\subsection{The legacy poset and its gaps}

The federation's credibility grades descend from a legacy five-element poset:
three rungs of evidential strength
(\emph{candidate} $<$ \emph{empirical} $<$ \emph{domain-formal}) plus two
cross-cutting grades (\emph{mirror-insight} for insight-track claims and
\emph{policy} for governance-track claims). Grades are used to gate which
claims may serve as premises in deduction, so the poset's \emph{join}
operation is load-bearing: the join of the grades of two premises is the
natural grade of any conclusion drawn from both.

An exhaustive enumeration over all pairs of the five-element poset
(certificate CERT-LATTICE-01) found that the poset has \emph{exactly seven
join gaps}: pairs with no least upper bound. The gaps are exactly the
cross-track pairs: each of the three rungs fails to join with
\emph{mirror-insight} and with \emph{policy} (six pairs), and
\emph{mirror-insight} fails to join with \emph{policy} (the seventh). A poset
with join gaps cannot support premise combination; completion is forced.

\subsection{The completion}

\begin{definition}[Credibility lattice]\label{def:lattice}
The carrier is $E = (R \times T) \cup \{\top, \bot\}$, where
$R = \{\text{candidate} < \text{empirical} < \text{domain-formal}\}$ is the
three-rung chain and
$T = \{\text{adjudication}, \text{insight}, \text{governance}\}$ is the set of
tracks. The order compares rungs within a track; $\bot$ is below everything
and $\top$ above everything; cells of different tracks are incomparable.
Join of two cells is the rung-maximum within a common track and $\top$
across tracks; meet is the rung-minimum within a common track and $\bot$
across tracks.
\end{definition}

\begin{remark}
The completion embeds the legacy poset order-preservingly and cleanly:
\emph{candidate}, \emph{empirical}, \emph{domain-formal} map to the
adjudication track at their rungs, \emph{mirror-insight} maps to
$(\text{candidate}, \text{insight})$, and \emph{policy} maps to
$(\text{candidate}, \text{governance})$. The enumeration reports
order-preservation true and reflection violations none (CERT-LATTICE-01).
\end{remark}

\begin{theorem}[CERT-LATTICE-01: exhaustive lattice verification]
\label{thm:lattice-enum}
Commutativity of join and meet, associativity of join and meet, and both
absorption laws hold for all $11^{3} = 1331$ ordered triples of the
eleven-element structure of Definition~\ref{def:lattice} (exhaustive machine
check, zero failures).
\end{theorem}

\subsection{Lean formalization}

The full structure---carrier as an inductive type, order, join, meet---is
formalized in Lean~4 (\lean{artifacts/lean-verified/CertLatticeLean01.lean},
certificate CERT-LATTICE-LEAN-01). Because the carrier is finite (eleven
elements), every universal statement over it is decidable, and every theorem
below is proved by the single tactic \texttt{decide}, which reduces to kernel
evaluation; cloud verification request
\rid{1dfa70b6-ba1e-492d-b45c-5391727aff3c} (Axiom Axle,
\texttt{lean-4.28.0}, okay\,=\,true).

\begin{theorem}[CERT-LATTICE-LEAN-01]\label{thm:lattice-lean}
Over the carrier of Definition~\ref{def:lattice}, all of the following hold
and are machine-decided in Lean~4:
\begin{enumerate}
\item the six equational lattice laws: \lean{join\_comm}, \lean{meet\_comm},
\lean{join\_assoc}, \lean{meet\_assoc}, and the two absorption laws
\lean{absorb1}, \lean{absorb2};
\item the order axioms: \lean{le\_refl}, \lean{le\_trans},
\lean{le\_antisymm};
\item extremality: \lean{join\_is\_lub} (the join is below every common upper
bound) and \lean{meet\_is\_glb} (the meet is above every common lower
bound);
\item order--operation compatibilities: \lean{le\_join\_left},
\lean{le\_join\_right}, \lean{meet\_le\_left}, \lean{meet\_le\_right}.
\end{enumerate}
\end{theorem}

\begin{remark}[Declaration count]
The certificate header describes a fourteen-theorem core package; the pinned
file contains fifteen \texttt{decide}-proved declarations in total, the
fifteenth being the fourth order--operation compatibility listed above. We
report the file's actual content.
\end{remark}

The axiom audit (Section~\ref{sec:repro}) confirms that these theorems depend
only on \texttt{propext}, \texttt{Classical.choice}, and
\texttt{Quot.sound}---the standard Lean~4 axiom whitelist---with no
\texttt{sorryAx}. Because the proofs are by \texttt{decide}, their content is
identical to the Python enumeration of Theorem~\ref{thm:lattice-enum}; the
value of the Lean layer is that the \emph{statement} is machine-checked for
well-formedness against a single source of truth, eliminating
transcription drift between the specification and the enumeration script.

\section{Certificate Lifecycle Machine}\label{sec:lifecycle}

\subsection{States and transitions}

Certificates in the federation are not static. A deposited certificate is
born as a \emph{candidate}, may be \emph{granted} after a four-gate review,
is \emph{maintained} under continuous monitoring, may be \emph{demoted} when
a domain violation is detected, and may be \emph{revoked} when forgery is
detected or its domain is withdrawn.

\begin{definition}[Lifecycle machine]\label{def:k4}
The state set is
$S=\{\text{candidate},\text{granted},\text{maintained},\text{demoted},\text{revoked}\}$.
The legality predicate $\ell : S \times S \to \{\mathrm{true},\mathrm{false}\}$
admits exactly the following nine transitions:
\begin{center}
\begin{tabular}{lll}
\toprule
From & To & Trigger \\
\midrule
candidate & granted & four-gate review \\
candidate & revoked & domain withdrawal \\
granted & maintained & continuous monitoring passed \\
granted & demoted & domain violation detected \\
granted & revoked & certificate forgery detected \\
maintained & demoted & domain violation detected \\
maintained & revoked & certificate forgery detected \\
demoted & granted & remediation re-review \\
demoted & revoked & forgery detected or domain withdrawal \\
\bottomrule
\end{tabular}
\end{center}
\end{definition}

\subsection{Invariants, including the one formalization broke}

\begin{theorem}[CERT-K4-01\,/\,CERT-K4-LEAN-01]\label{thm:k4}
The machine of Definition~\ref{def:k4} satisfies, all by exhaustive machine
check and all re-proved in Lean~4 by \texttt{decide} over the five-state
carrier (cloud verification request
\rid{16618831-e457-4d29-8d65-b57fa563e33d}):
\begin{enumerate}
\item \lean{legal\_count}: exactly $9$ transitions are legal;
\item \lean{illegal\_count}: exactly $11$ distinct (non-self) transitions are
illegal;
\item \lean{I1\_granted\_gated} (grant gating): if $\ell(s,\text{granted})$
then $s\in\{\text{candidate},\text{demoted}\}$;
\item \lean{I1b\_maintained\_only\_from\_granted}: if
$\ell(s,\text{maintained})$ then $s=\text{granted}$;
\item \lean{I2\_revoked\_terminal} (revocation terminality): no transition
leaves \emph{revoked};
\item \lean{I2b\_no\_return\_to\_candidate} (irreversibility): no transition
returns to \emph{candidate};
\item \lean{I3\_demoted\_out\_edges} (appeal freezing): if
$\ell(\text{demoted},s)$ then $s\in\{\text{granted},\text{revoked}\}$;
\item \lean{no\_self\_loop}: $\ell(s,s)$ is false for every $s$.
\end{enumerate}
\end{theorem}

\begin{remark}[A specification flaw caught by formalization]\label{rem:i1}
Our original specification stated invariant (I1) more strongly: ``\emph{granted}
arises only from \emph{candidate}''. This is \emph{false}: the
remediation edge $\text{demoted}\to\text{granted}$ is a legal entry into
\emph{granted}. The Lean formalization's \texttt{decide} procedure produced a
counterexample to the original statement, and the invariant was corrected to
the disjunctive form in Theorem~\ref{thm:k4}(3), which is the machine-proved
one. We regard this incident as the methodological point of the exercise:
the informal specification, reviewed by humans and by machines at the
enumeration layer, had silently dropped an edge from the gating invariant;
only the act of stating the invariant in a proof assistant exposed the gap.
The corrected statement is recorded as the ``Lean-ification precision
correction'' in the certificate header of
\lean{CertK4Lean01.lean}.
\end{remark}

The Python enumeration layer (CERT-K4-01, pinned with raw data in the
federation ledger) independently confirms the counts and the three invariants
over the full transition graph, including the enumeration of all eight
candidate-originating paths for the append-only evidence property and the
closure of all five states' outgoing edges under an raised appeal flag.

\section{Certified Numerical Anchors}\label{sec:anchors}

The kernel's adjudication machinery would be vacuous without at least one
class of certificates whose acceptance is anchored in mathematics rather than
in trust. This section presents two such anchor families and then states,
as a first-class finding, the boundary of the Lean-side port.

\subsection{Circulant Sinkhorn fixed points: the Krawczyk family}
\label{sec:circ}

For a circulant cost matrix $C[i][j]=c_{(j-i)\bmod k}$ with
$c=(c_0,\dots,c_{k-1})\in\mathbb{Q}_{>0}^{k}$, uniform marginals,
regularization $\varepsilon>0$, and the gauge $f_0=0$, the entropic
optimal-transport (Sinkhorn) fixed-point system reads
\[
f_i+\varepsilon\log k+\varepsilon\log\sum_{j} e^{(g_j-C[i][j])/\varepsilon}=0,
\qquad
g_j+\varepsilon\log k+\varepsilon\log\sum_{i} e^{(f_i-C[i][j])/\varepsilon}=0 .
\]
Circulant symmetry forces the closed-form solution
\[
f^{*}=0,\qquad g^{*}=-\varepsilon\bigl(\log k+\operatorname{lse}(-c/\varepsilon)\bigr),
\]
where $\operatorname{lse}$ is the log-sum-exp; e.g.\ for the $k=6$,
$\varepsilon=1$ instance with $c=(1/2,1/3,2/3,1/4,3/4,1/5)$, $g^{*}=0$ in
the gauge and the common $g$-value is $-3.1543698608222522\ldots$, reproduced
digit-for-digit by the independent batch verification. The closed form is
used only as ground truth for negative controls; the certificate does not
consult it.

The certificate is the classical Krawczyk criterion~\cite{krawczyk1969}. With
$F$ the $(2k-1)$-dimensional gauged system, $\hat{x}$ a rationalized center
(rational approximation of the closed form, $\texttt{limit\_denominator}$ at
precision $10^{-13}$ for the original batch and $\sim 10^{-26}$ for the
re-verification batch), $C$ a rationalization of the analytic inverse
Jacobian $J(\hat{x})^{-1}$, and $X$ the box of radius $10^{-9}$ around
$\hat{x}$, the operator
\[
K(X) = \hat{x} - C\,F(\hat{x}) + \bigl(I - C\,J(X)\bigr)(X-\hat{x})
\]
is evaluated in interval arithmetic with outward rounding (mpmath interval
layer, 80 decimal digits). If $K(X)$ lies strictly inside $X$ and $C$ is
nonsingular, $X$ contains exactly one real root.

\begin{theorem}[CERT-CIRC-01 family; Python interval certification]
\label{thm:circ}
For each of the six instances $k\in\{6,10\}$,
$\varepsilon\in\{1,\tfrac12,\tfrac15\}$ (two circulant vector families,
dimensions $n=11$ and $n=19$), the Krawczyk operator maps the box strictly
inside itself: Jacobian condition numbers $22.2$--$43.1$, Krawczyk image
widths between $2.55\cdot10^{-15}$ and $6.93\cdot10^{-14}$, against a box
radius of $10^{-9}$ (margin of at least five orders of magnitude). Hence each
box contains exactly one real root. Every negative control---the same box
with centers perturbed by $10^{-6}$, preconditioner recomputed---is rejected.
An independent re-verification batch (EXT05-BATCH-01) reproduced all six
instances (condition numbers matching to published precision) and extended
the family to $\varepsilon\in\{1/10,1/20\}$, for a total of $10/10$
instances with strict inclusion and $10/10$ negative controls correctly
rejected; at $\varepsilon=1/20$ the condition number rises to
$1.1\cdot10^{3}$ ($k=6$) and the image width stays below $5\cdot10^{-14}$.
\end{theorem}

\begin{remark}
The extension to small $\varepsilon$ is where the certification earns its
keep: the closed form is cheap, but the \emph{box uniqueness} statement is
what downstream adjudication consumes, and it survives condition numbers
three orders of magnitude above the original family.
\end{remark}

\subsection{Auction-algorithm assignment: the dual-bracket family}
\label{sec:auction}

The second anchor family (CERT-AUCTION-01) certifies the output of the
auction algorithm~\cite{bertsekas1988} on assignment problems with
log-uniform costs $C=10^{R\,(2U-1)}$, $U$ uniform on $[0,1]^{k\times k}$,
dynamic range $R=2$. The solver (forward auction with
$\varepsilon$-scaling, $\varepsilon: 0.5\to/5\to 10^{-6}$ per round, warm
prices) is treated as an untrusted candidate generator; the certificate is a
posteriori and rests only on LP weak duality: an interval evaluation (outward
rounding via directed rounding of endpoints) certifies reduced-cost
nonnegativity $\mathrm{rc}_j = c_j - (A^{T}y)_j \ge 0$ componentwise, with a
$\delta$-deflation ($\delta=10^{-10}$) at degenerate zero-riding coordinates;
the assignment-scale bracket is $k\times[\mathrm{LB}_{\mathrm{lo}},
\mathrm{UB}_{\mathrm{hi}}]$ in the Birkhoff scaling.

\begin{proposition}[CERT-AUCTION-01; Python interval certification]
\label{prop:auction}
All five instances (one benchmark at $k=8$, $R=2$, seed $7$; three extended
instances at $k\in\{6,10,12\}$; one schedule-variant re-run of the
benchmark) pass the three-part test: the auction cost lies in the certified
bracket, the gap satisfies $\mathrm{cost}-k\cdot\mathrm{LB}_{\mathrm{lo}}\le
k\,\varepsilon_{\mathrm{final}}$, and a swap-perturbation negative control
(perturbed cost $97.28$--$130.71$ against brackets of width
$\sim 5\cdot10^{-11}$) falls outside the bracket and is rejected. The
benchmark reproduces the originally registered cost
$0.2550220110027003$ digit-for-digit, and the schedule-variant converges to
the same assignment, family-level evidence of path independence of
$\varepsilon$-scaling.
\end{proposition}

\begin{remark}[Honesty boundary of the auction family]\label{rem:auction-honest}
The upper side of the bracket uses the LP solver's primal solution as the
candidate; primal feasibility is trusted to solver precision
($\sim 10^{-12}$) and is \emph{not} interval-certified---this is the original
specification, and we flag it rather than repair it silently. The family has
no Lean port.
\end{remark}

\subsection{The boundary of the Lean port (FINDING-04)}\label{sec:finding04}

We attempted to port the six circulant instances to Lean~4 using the verified
interval library \emph{leancert}, encoding each system over \emph{leancert}'s
expression grammar and calling
\lean{verify\_unique\_system\_root} with \texttt{native\_decide}. The port
does not go through, and the failure is structural, not incidental:

\begin{enumerate}
\item \emph{leancert}'s supported-expression check
(\lean{ADSupported}\,/\,\lean{checkADSupported}) does not include
\lean{Expr.log}---it never has, since the commit that introduced the
function---while all $11/11$ (resp.\ $19/19$) equations of every circulant
system contain $\log$. The checker therefore rejects by design: probes
measure \lean{checkADSupported}\,=\,false on $11/11$ equations while
\lean{centerInside}\,=\,true.
\item \emph{leancert}'s determinant (\lean{Matrix.det}) is defined by
Leibniz expansion; at dimensions $11$ and $19$ the computation is
combinatorially infeasible for \texttt{native\_decide} (the library's own
tests cover only $2\times 2$).
\end{enumerate}

Consequently, under the artifact-pinned \emph{leancert} v4.34.1 the
circulant Lean family \emph{has never passed} Lean-side machine verification:
after elaboration-layer repair the files elaborate cleanly
(\textsc{elaboration ok}), but \texttt{native\_decide} evaluates to false for
the structural reasons above, and the repaired artifacts carry the honest
verdict \textsc{check\_false (structural)}. \textbf{The certification of the
circulant family is carried entirely by the Python interval layer}
(outward rounding, strict-inclusion criterion, nonsingular preconditioner,
negative controls), and the Lean files are \emph{formal statement
artifacts}---machine-checked transcriptions of the systems, centers, boxes,
and claims---not machine-verified certificates. A content-level repair path
is registered as an open obligation: eliminate the logarithm by the
substitution $u=\log\sum\exp(\cdot)$ to bring the systems into the
library's supported fragment, or migrate to a checker with transcendental
support; the rationalized centers for the extended instances are already
registered for this purpose.

We document this boundary in the body of the paper, rather than in a
footnote, because it is exactly the kind of fact that an adjudication kernel
exists to protect: a credential whose presentation outruns its certification
layer. The v1 abstract's wording (``ported to \emph{leancert} with
\lean{verify\_unique\_system\_root} certificates'') was, under this finding,
overstated, and is corrected here.

\section{A Parameterized Rice Template}\label{sec:rice}

\subsection{Bridge, extensionality, generator}

The domain-limitation architecture of Section~\ref{sec:intro} is justified by
Rice's theorem~\cite{rice1953}, and Mathlib contains a formalization of that
theorem over partial-recursive-function codes~\cite{carneiro2019}. The
kernel's fourth component (T2A-BRIDGE-01, CERT-T2A-TEMPLATE-01) packages the
blocking result as a callable template in Lean~4. All six theorems below were
cloud-verified under \texttt{verify\_proof} (Axiom Axle,
\texttt{lean-4.28.0}, okay\,=\,true, request identifiers in
Table~\ref{tab:evidence}), with axiom audits clean.

\begin{proposition}[Bridge, \lean{rice\_bridge}]\label{prop:bridge}
Let $C\subseteq \mathrm{Code}$ be extensionally closed (codes with equal
evaluation semantics are equimembership) and nontrivial
($C\neq\emptyset$, $C\neq\mathrm{univ}$). Then membership in $C$ is not a
computable predicate.
\end{proposition}

The proof is a three-line reduction to Mathlib's \lean{ComputablePred.rice\_2}:
decidability of $C$ would force $C=\emptyset$ or $C=\mathrm{univ}$.

\begin{lemma}[Extensionality of pointwise predicates, \lean{ext\_of\_pointwise}]
\label{lem:ext}
For any $g:\mathbb{N}\to\mathbb{N}$, the predicate
``$\exists m,\ \mathrm{eval}\ c\ m = \mathrm{some}\ (g\,m)$'' is
extensionally closed in the sense of Proposition~\ref{prop:bridge}.
\end{lemma}

\begin{theorem}[Generator, \lean{rice\_pointwise}]\label{thm:gen}
For any $g:\mathbb{N}\to\mathbb{N}$ whose range avoids some point $k$
(i.e.\ $\forall m,\ g\,m\neq k$), the pointwise predicate of
Lemma~\ref{lem:ext} is undecidable.
\end{theorem}

\begin{proof}[Proof sketch]
Instantiate the bridge with $C=\{c \mid \exists m,\ \mathrm{eval}\ c\
m=\mathrm{some}\ (g\,m)\}$. Extensionality is Lemma~\ref{lem:ext}.
Nontriviality is witnessed by the constant functions: $\mathrm{const}\ (g\,0)$
is in $C$ (at $m=0$), and $\mathrm{const}\ k$ is not (the range-avoidance
hypothesis). The Lean proof is a direct formalization of this argument.
\end{proof}

\subsection{Discharged instances}

The generator is discharged on seven instances of increasing witness
difficulty (Table~\ref{tab:rice}). The first three accompany the template
file itself; the remaining four (INSTANCES-EXT01) demonstrate that the
generator's call surface extends from constant and linear witnesses to
nonlinear arithmetic, inductive witnesses, and divisibility arguments. All
seven carry cloud verification request identifiers and clean axiom audits;
the EXT01 batch was additionally verified as a single four-theorem artifact
(request \rid{77a0eaad-fbd2-48a6-9281-330f736fc8c2}).

\begin{table}[t]
\centering
\caption{Discharged instances of the Rice template.
``Avoids'' is the point $k$ of Theorem~\ref{thm:gen}.}\label{tab:rice}
\begin{tabular}{llll}
\toprule
Instance & Semantic property ($g$) & Avoids & Witness content \\
\midrule
\lean{inst\_const\_zero} & output constantly $0$ & $k=1$ & \lean{Nat.zero\_ne\_one} \\
\lean{inst\_succ} & output $=$ input $+1$ & $k=0$ & \lean{Nat.succ\_ne\_zero} \\
\lean{inst\_double} & output $=2m$ & $k=1$ & parity (\texttt{omega}) \\
\lean{inst\_shift7} (E4) & output $=$ input $+7$ & $k=3$ & $\beta$-redex elimination + \texttt{omega} \\
\lean{inst\_square\_ne3} (E5) & output $=m^{2}$ & $k=3$ & case split: $m<2$ enumeration; $m\ge2\Rightarrow m^{2}\ge4$ \\
\lean{inst\_pow2\_ne3} (E6) & output $=2^{m}$ & $k=3$ & induction: $2^{0}=1$, $2^{m+1}$ even \\
\lean{inst\_fact\_ne4} (E7) & output $=m!$ & $k=4$ & divisibility: $m\ge3\Rightarrow 6\mid m!$ \\
\bottomrule
\end{tabular}
\end{table}

\begin{remark}[Engineering note]\label{rem:beta}
The EXT01 batch recorded a reusable engineering fact: Lean's \texttt{omega}
tactic does not see through $\beta$-redexes, so instance proofs first reduce
the goal with \texttt{show} before invoking \texttt{omega} or \texttt{rw}.
This is the kind of artifact-level detail that the federation ledger retains
and that a template's reusability depends on.
\end{remark}

\section{Minimal Trusted Base}\label{sec:selfcheck}

Every anchor certificate ultimately reduces to: a checker computed
\emph{accept}, therefore the semantic property holds. The kernel's fifth
component (CERT-SELFCHECK-01) proves this pattern, in miniature, in Lean~4,
over rational intervals.

\begin{definition}[$\mathbb{Q}$-intervals and the checker]
An interval \lean{Iv} is a pair $(\mathrm{lo},\mathrm{hi})\in\mathbb{Q}^{2}$;
membership \lean{mem} is $\mathrm{lo}\le x\le\mathrm{hi}$. Interval addition
and negation are endpointwise. The boundedness checker
\lean{check}\ $(\mathrm{iv},B)$ is the proposition $-B\le\mathrm{lo}\ \wedge\
\mathrm{hi}\le B$.
\end{definition}

\begin{theorem}[CERT-SELFCHECK-01: accept implies correctness]
\label{thm:selfcheck}
The following are proved in Lean~4 (cloud verification under
\texttt{verify\_proof}, \texttt{lean-4.28.0}, okay\,=\,true; request
identifiers in Table~\ref{tab:evidence}; axiom audit clean):
\begin{enumerate}
\item \lean{add\_sound}: interval addition preserves membership;
\item \lean{neg\_sound}: interval negation preserves membership;
\item \lean{check\_sound}: if \lean{check}\ $(\mathrm{iv},B)$ holds, then
every $x\in\mathrm{iv}$ satisfies $|x|\le B$;
\item \lean{end\_to\_end}: composition---if the certificate arithmetic
$\mathrm{add}\ a\ (\mathrm{neg}\ b)$ passes \lean{check} against $B$, then
$|x-y|\le B$ for all $x\in a$, $y\in b$.
\end{enumerate}
\end{theorem}

\begin{remark}
This is deliberately tiny. Its role in the kernel is foundational rather
than mathematical: it is a machine-checked instance of the soundness
interface $V(\varphi,D,\pi)=\mathrm{accept}\Rightarrow\varphi$, exhibiting
that the interface is satisfiable by a real, verified artifact. The full
Krawczyk-layer analogue---a verified interval evaluator with outward rounding
strong enough to carry Theorem~\ref{thm:circ} inside Lean---is precisely
what FINDING-04 (Section~\ref{sec:finding04}) shows to be beyond the pinned
\emph{leancert} version; it is registered as open obligation OBL-A1.
\end{remark}

\section{Reproducibility and Axiom Audit}\label{sec:repro}

\subsection{Pinned artifacts and verification channels}

All artifacts are pinned by path and commit in the public repository
\texttt{chepin-ai/vci-inbox}; the index is in Appendix~\ref{app:anchors}.
The verification channels are of three kinds:
\begin{itemize}
\item \emph{Exhaustive enumeration} (Python, pinned scripts): lattice laws
over $1331$ triples; lifecycle transition counts and invariants. Raw data
are frozen in the ledger (board/LAB-THEORY-01R-FK01R-20261009T0620Z.md,
commit \rid{3e0f54e1}; reproduction script commit \rid{16ed41d2}).
\item \emph{Interval certification} (Python, mpmath interval layer,
outward rounding, 80-digit precision): circulant Krawczyk family
(EXT05-BATCH-01 script and machine-readable results under
\lean{artifacts/CERT-CIRC-LEAN-01/family/}) and the auction dual-bracket
family (\lean{artifacts/CERT-AUCTION-01/auction\_verify.py}).
\item \emph{Cloud proof verification} (Axiom Axle, environment
\texttt{lean-4.28.0}, strict \texttt{verify\_proof} channel: no
\texttt{sorry}, whitelisted axioms, signature match): all Lean artifacts of
Sections~\ref{sec:lattice}--\ref{sec:selfcheck}. Each theorem carries a
permanent request identifier under which the verification can be re-queried;
the identifiers are listed in Table~\ref{tab:evidence} and in the artifact
headers.
\end{itemize}

\subsection{Axiom audit}

Every Lean artifact was audited in the cloud environment with
\texttt{\#print axioms}. In every case the dependency set is exactly a subset
of the standard Lean~4 whitelist
\[
\{\texttt{propext},\ \texttt{Classical.choice},\ \texttt{Quot.sound}\},
\]
and \texttt{sorryAx} appears in none. Because the lattice and lifecycle
theorems are proved by \texttt{decide}, their logical content reduces to
kernel computation; the choice axiom enters only through Mathlib's
\texttt{Finset} machinery and is inessential to the checked statements.

\subsection{Environment fingerprint and dependency-drift caveat}

The cloud verification environment fingerprint is \texttt{lean-4.28.0}
(Axiom Axle). The attempted \emph{leancert} port additionally records a
dependency-pinning incident relevant to reproducibility: the upstream tags
v4.34.1\,/\,v4.35.0 of \emph{leancert} were observed to have been retargeted
on the remote (the tag target differs from the artifact-pinned commit), so
the artifact pins by commit hash and the drift risk is registered in the
federation ledger. We reproduce this caveat because any consumer of the
formal statement artifacts must resolve the dependency by the pinned commit,
not by tag.

\subsection{Empirical soundness: fuzzing with rule-of-three bounds}

Three certificate families were fuzzed (certificate CERT-T4-01, raw data in
the pinned ledger card; $f_{80}$ ground truth means containment of the
80-bit-extended-precision value, judged bitwise):
\begin{itemize}
\item \textbf{E family} (interval envelopes): $n=30$, ground-truth
containment $30/30$;
\item \textbf{D family} (LP dual brackets, with implanted corrupted
subsets): $n=60$, accepts $4$, all $4$ valid;
\item \textbf{K family} (Krawczyk, adversarial centers with implanted true
center): $n=30$, accept $1$, independently verified by an $f_{80}$ Newton
iteration.
\end{itemize}
Zero false accepts were observed. The three layers are deliberately
non-identically distributed and we do not extrapolate a universal statement.
By the rule of three~\cite{hanley1983}, the 95\% upper confidence bounds on
the false-accept rate are $9.5\%$ (E), $4.9\%$ (D), $9.5\%$ (K).
\textbf{We report ``no false accept observed'', not a guarantee.}

\subsection{Evidence-anchor table}

\begin{table}[t]
\centering
\caption{Evidence anchors. rid = cloud verification request identifier
(Axiom Axle \texttt{verify\_proof}, \texttt{lean-4.28.0}, okay\,=\,true,
axioms within \{\texttt{propext}, \texttt{Classical.choice},
\texttt{Quot.sound}\}). Truncated rids are abbreviated with $\ldots$; full
identifiers appear in the pinned artifact headers.}\label{tab:evidence}
\small
\begin{tabular}{>{\raggedright\arraybackslash}p{4.6cm} >{\raggedright\arraybackslash}p{4.4cm} >{\raggedright\arraybackslash}p{5.2cm}}
\toprule
Artifact & Verification status & Anchors \\
\midrule
CERT-LATTICE-01 (enum.) & exhaustive, $1331$ triples, $0$ failures &
commit \rid{3e0f54e1}; repro \rid{16ed41d2} \\
CERT-LATTICE-LEAN-01 & $15$ theorems by \texttt{decide}; audit clean &
rid \rid{1dfa70b6-\ldots}; commit \rid{3a5edd44} \\
CERT-K4-01 (enum.) & $9$ legal, $11$ illegal, I1--I3 pass &
commit \rid{3e0f54e1} \\
CERT-K4-LEAN-01 & $8$ theorems by \texttt{decide}; audit clean; I1 corrected &
rid \rid{16618831-\ldots}; commit \rid{3a5edd44} \\
CERT-CIRC-01 family ($6$ inst.) & Python interval: strict inclusion $6/6$,
neg.\ controls reject $6/6$ & \lean{artifacts/CERT-CIRC-LEAN-01}, commits
\rid{7d052564}, \rid{bd71b720} \\
EXT05-BATCH-01 ($10$ inst.) & Python interval: $10/10$ inside, $10/10$
neg.\ controls & \lean{family/EXT05-BATCH-01\_verify.py}, results JSON \\
CERT-CIRC-LEAN-01 (Lean) & formal statement only;
\textsc{elaboration ok}, \textsc{check\_false (structural)} (FINDING-04) &
\lean{*\_v4341.lean}, commit \rid{c8c43b39} \\
CERT-AUCTION-01 ($5$ inst.) & Python interval: $5/5$ pass; neg.\ controls
reject; primal feasibility not interval-certified (Remark~\ref{rem:auction-honest}) &
\lean{auction\_verify.py}, results JSON \\
T2A-BRIDGE-01 & \lean{t2a\_bridge}, \lean{t2a\_instance} verified &
rids \rid{c2a58cfa-\ldots}, \rid{a446bfd2-\ldots}; commit \rid{ebc121a1} \\
CERT-T2A-TEMPLATE-01 & $6$ theorems verified; audit clean &
rids \rid{41e07431-\ldots}--\rid{2e3b960c-\ldots}; commit \rid{577b1a4f} \\
INSTANCES-EXT01 (E4--E7) & $4$ instances verified; batch rid
\rid{77a0eaad-\ldots} & \lean{artifacts/T2A-BRIDGE-01/INSTANCES-EXT01.lean} \\
CERT-SELFCHECK-01 & $4$ theorems verified; audit clean &
rids \rid{d9034a05-\ldots}, \rid{44208d46-\ldots}, \rid{082321e6-\ldots},
\rid{f3c1fc61-\ldots}; commit \rid{f8cb83e7} \\
CERT-T4-01 (fuzz) & $91$ trials, $0$ false accepts observed;
rule-of-three bounds reported & commit \rid{3e0f54e1} \\
\bottomrule
\end{tabular}
\end{table}

\section{Limitations and Honest Gaps}\label{sec:limitations}

We collect the limitations in one place, matching the federation's honesty
discipline that every claim be traceable to an artifact and every gap be
named rather than elided.

\begin{enumerate}
\item \textbf{Machine generation and human comprehension.}
This paper and all of its artifacts are machine-generated under the
authorship mode \emph{no-human-author-claimed}. The contributing account
holder does not claim human understanding of every derivation step; the
correctness belief rests on the verification layers enumerated in
Section~\ref{sec:repro}, not on human line-by-line review. Human
comprehension of some derivation steps is therefore partial by design, and
readers should weight the machine certificates accordingly---this cuts both
ways: the certificates are stronger than a human read in exhaustiveness, and
weaker in semantic judgment about what \emph{should} have been proved.
\item \textbf{The circulant Lean boundary.} The circulant family's
certification is carried by the Python interval layer. The Lean-side files
are formal statement artifacts that elaborate but cannot be machine-verified
under the pinned \emph{leancert} v4.34.1, for the two structural reasons of
Section~\ref{sec:finding04} (no \lean{Expr.log} support; infeasible Leibniz
determinant at $n\in\{11,19\}$). Any v1-era wording implying a completed
Lean-side certificate for this family was overstated and is corrected here.
The repair path ($\log$-elimination by substitution, or a checker with
transcendental support) is an open obligation, not a result.
\item \textbf{The auction family's trusted edge.} The dual bracket's upper
side trusts solver primal feasibility to $\sim 10^{-12}$; it is not
interval-certified (Remark~\ref{rem:auction-honest}). The family has no Lean
port.
\item \textbf{No statistical guarantees.} All fuzzing conclusions are
sample-based, non-identically distributed across layers, and reported with
rule-of-three upper bounds. ``Zero false accepts observed'' is the complete
claim.
\item \textbf{Finite and small carriers.} The lattice and lifecycle results
are decidable finite checks. They are complete for the stated structures, but
they do not by themselves establish properties of any larger or dynamical
system built on top of them.
\item \textbf{Trusted bases we do not eliminate.} The kernel assumes a
finite-depth checking stack terminating in hardware, microcode, and audit
anchors; it does not claim an ultimate foundation. The checker-reliability
interface at the Krawczyk and LP layers is discharged by classical theorems
(outward-rounding inclusion, Krawczyk's theorem, weak duality), not yet by
proof-assistant formalization at full strength.
\item \textbf{Open theses are not theorems.} The federation's escape-route
catalogue around Rice-type blockage (domain limitation, probabilistic
checking, interactive proofs, restricted fragments, paraconsistent or
many-valued semantics, semi-decision) is registered as an open, falsifiable
thesis with an explicit retraction of earlier universal wordings; none of it
is claimed in this paper.
\end{enumerate}

\section{Conclusion}\label{sec:conclusion}

We have deposited, in full paper form, a machine-checked kernel for
federated automated adjudication: a completed credibility lattice verified
by exhaustive enumeration and by Lean~4 decision over the finite carrier; a
certificate lifecycle machine whose formalization caught and corrected a real
specification flaw; two certified numerical anchor families, with an exact
and honest account of which layer carries the certification and where the
Lean port is structurally blocked; a parameterized Rice template with seven
discharged instances making the undecidability boundary callable; and a
minimal, machine-proved trusted base for the accept-implies-correct pattern.
Every claim is pinned to an artifact, every Lean theorem to a
cloud-verification request identifier, every axiom audit to the standard
whitelist, and every empirical statement to a rule-of-three bound. The
project's broader question---whether a federation of machines can maintain
adjudication infrastructure whose credentials never outrun their
certificates---is answered here only at kernel scale, but the kernel is
complete, pinned, and falsifiable in every part.

\appendix

\section{Pinned Artifact Index}\label{app:anchors}

All paths are relative to the repository root of \texttt{chepin-ai/vci-inbox}
(branch \texttt{main}); commit identifiers pin the introducing or repairing
commit as recorded in the federation ledger.

\begin{itemize}
\item \lean{artifacts/lean-verified/CertLatticeLean01.lean} (commit
\rid{3a5edd44}): lattice formalization, Section~\ref{sec:lattice}.
\item \lean{artifacts/lean-verified/CertK4Lean01.lean} (commit
\rid{3a5edd44}): lifecycle formalization, Section~\ref{sec:lifecycle}.
\item \lean{artifacts/CERT-CIRC-LEAN-01/} (commits \rid{7d052564},
\rid{bd71b720}; repaired \lean{*\_v4341.lean} at \rid{c8c43b39}): circulant
family, Sections~\ref{sec:circ} and~\ref{sec:finding04}; the
\lean{family/} subdirectory holds the batch re-verification script, results,
and the extended instances.
\item \lean{artifacts/CERT-AUCTION-01/}: auction family,
Section~\ref{sec:auction} (\lean{auction\_verify.py}, results JSON, README
with the honesty boundary).
\item \lean{artifacts/T2A-BRIDGE-01/t2a\_bridge.lean} (commit
\rid{ebc121a1}) and \lean{artifacts/lean-verified/T2aTemplate.lean} (commit
\rid{577b1a4f}); \lean{artifacts/T2A-BRIDGE-01/INSTANCES-EXT01.lean}:
Rice template and instances, Section~\ref{sec:rice}.
\item \lean{artifacts/lean-verified/CertSelfCheck01.lean} (commit
\rid{f8cb83e7}): minimal trusted base, Section~\ref{sec:selfcheck}.
\item \lean{board/LAB-THEORY-01R-FK01R-20261009T0620Z.md} (commit
\rid{3e0f54e1}; reproduction script commit \rid{16ed41d2}): enumeration
certificates CERT-LATTICE-01, CERT-K3-01, CERT-K4-01, CERT-T4-01 with raw
data, Section~\ref{sec:repro}.
\item \lean{board/FINDING-20261010-04-*.md}: the structural boundary of the
circulant Lean port, Section~\ref{sec:finding04}.
\end{itemize}

\section{Cloud Verification Request Identifiers}\label{app:rids}

For completeness, the full request identifiers recorded in the artifact
headers (Axiom Axle \texttt{verify\_proof}, environment
\texttt{lean-4.28.0}, all okay\,=\,true):
\begin{itemize}
\item Lattice: \rid{1dfa70b6-ba1e-492d-b45c-5391727aff3c}.
\item Lifecycle: \rid{16618831-e457-4d29-8d65-b57fa563e33d}.
\item Rice bridge (T2A-BRIDGE-01): \rid{c2a58cfa-dd32-4bcb-a8c1-b06f93d25c08},
\rid{a446bfd2-d9a0-4727-afe3-d5f13b63a401}.
\item Rice template: \rid{41e07431-e53b-465d-b920-c4898fa257c0},
\rid{2d611c78-046b-4e1f-91bf-7ab9a2a59aa3},
\rid{c39c5b84-bc93-4415-a5e2-fda74020a265},
\rid{4ee3fbb4-a0e0-4319-a270-67df7e883082},
\rid{d6fba622-b277-4e67-8927-da561d8ab989},
\rid{2e3b960c-7916-4f11-a608-43da3d4f74ca}.
\item Rice EXT01 instances: \rid{c5674825-412d-4877-b5cf-f8ae967ffe87},
\rid{bd879da4-b6c1-476c-b490-5d389374d438},
\rid{fdedfb82-095d-45dc-90ec-688e83bbdd6b},
\rid{17c045dd-cefe-480d-a70c-995696508739}; batch
\rid{77a0eaad-fbd2-48a6-9281-330f736fc8c2}.
\item Self-check: \rid{d9034a05-1956-43e9-87b0-c2bfe13ad989},
\rid{44208d46-9489-4670-b171-91099fa33ad9},
\rid{082321e6-6aeb-4a89-b51c-ac2ab5e926e1},
\rid{f3c1fc61-c0c0-4972-920e-9a3856c6453d}.
\end{itemize}

\begin{thebibliography}{99}

\bibitem{lean4}
L.~de Moura and S.~Ullrich.
\newblock The Lean~4 theorem prover and programming language.
\newblock In \emph{Automated Deduction -- CADE-28}, Lecture Notes in Computer
Science 12699, pp.~625--635. Springer, 2021.

\bibitem{mathlib}
The mathlib Community.
\newblock The Lean mathematical library.
\newblock In \emph{Proceedings of the 9th ACM SIGPLAN International Conference
on Certified Programs and Proofs (CPP 2020)}, pp.~367--381. ACM, 2020.

\bibitem{carneiro2019}
M.~Carneiro.
\newblock Formalizing computability theory via partial recursive functions.
\newblock In \emph{10th International Conference on Interactive Theorem
Proving (ITP 2019)}, LIPIcs 141, pp.~12:1--12:17. Schloss
Dagstuhl, 2019.

\bibitem{lean4lean}
M.~Carneiro.
\newblock Lean4Lean: Towards a formalized metatheory for the Lean theorem
prover.
\newblock arXiv:2403.14064, 2024.

\bibitem{rice1953}
H.~G. Rice.
\newblock Classes of recursively enumerable sets and their decision problems.
\newblock \emph{Transactions of the American Mathematical Society} 74(2):
358--366, 1953.

\bibitem{moore1966}
R.~E. Moore.
\newblock \emph{Interval Analysis}.
\newblock Prentice-Hall, Englewood Cliffs, NJ, 1966.

\bibitem{krawczyk1969}
R.~Krawczyk.
\newblock Newton-Algorithmen zur Bestimmung von Nullstellen mit
Fehlerschranken.
\newblock \emph{Computing} 4:187--201, 1969.

\bibitem{sinkhorn1967}
R.~Sinkhorn.
\newblock Diagonal equivalence to matrices with prescribed row and column
sums.
\newblock \emph{The American Mathematical Monthly} 74(4):402--405, 1967.

\bibitem{cuturi2013}
M.~Cuturi.
\newblock Sinkhorn distances: Lightspeed computation of optimal transport.
\newblock In \emph{Advances in Neural Information Processing Systems 26
(NIPS 2013)}, pp.~2292--2300, 2013.

\bibitem{bertsekas1988}
D.~P. Bertsekas.
\newblock The auction algorithm: A distributed relaxation method for the
assignment problem.
\newblock \emph{Annals of Operations Research} 14:105--123, 1988.

\bibitem{hanley1983}
J.~A. Hanley and A.~Lippman-Hand.
\newblock If nothing goes wrong, is everything all right? Interpreting zero
numerators.
\newblock \emph{Journal of the American Medical Association} 249(13):
1743--1745, 1983.

\end{thebibliography}

\end{document}
